	\documentclass[12pt]{exam}
	\usepackage[utf8]{inputenc}
	\usepackage[portuguese,brazil]{babel}
	\usepackage[margin=.65in]{geometry}
	\usepackage{amsmath,amssymb}
	\usepackage{multicol}
	\usepackage{multirow}
	\usepackage{float}
	\usepackage{enumerate}
	\usepackage{enumitem}
	\usepackage{array}
	\usepackage{graphicx}
	\usepackage{listings}
	\usepackage{circuitikz}    
	\usepackage{xcolor}
	\usepackage{tikz}
	\usetikzlibrary{shapes.multipart, calc}
	\usepackage{rotating}
	\usepackage{caption}
	%\usepackage[version=3]{mhchem}

	\pointpoints{ponto}{pontos}
	\hpword{Pontos:}
	\vpword{Pontos}
	\htword{Total}
	\vtword{Total:}
	\vsword{Resultado}

	\definecolor{azul}{RGB}{0,32,96}
	\definecolor{verde}{RGB}{0,132,124}
	\definecolor{cinza}{RGB}{0,84,132}
	%\definecolor{cinza2}{RGB}{218,220,216}
	\definecolor{cinza2}{RGB}{132,134,136}
	\definecolor{mygreen}{rgb}{0,0.6,0}
	\definecolor{mygray}{rgb}{0.5,0.5,0.5}
	\definecolor{mymauve}{rgb}{0.58,0,0.82}
	\newcommand*{\vttfamily}{%
	\fontencoding{T1}\fontfamily{lmvtt}\selectfont
	}
	\lstdefinestyle{Cstyle}{%
	language=C,numbers=left, keywordstyle=\color{mymauve}\vttfamily,commentstyle=\color{mygreen}\vttfamily,stringstyle=\color{blue},basicstyle=\scriptsize\vttfamily,numberstyle=\tiny\vttfamily,breaklines=true,stepnumber=1,tabsize=2,xleftmargin=\leftmarginii,escapechar=!,captionpos=b,morekeywords={include, printf,NULL},columns=flexible,escapeinside={(^}{^)},
	frame=l,
	framesep=2.5mm,
	framexleftmargin=2.5mm,
	fillcolor=\color{cinza2!50},
	rulecolor=\color{ballblue},
	rulecolor=\color{blue}
	}

	\lstnewenvironment{C}[1][]
	{
	\lstset{style=Cstyle,#1}
	}
	{}


	\newcommand{\class}{Segurança da Informação}
	\newcommand{\profname}{Gilmar Gomes do Nascimento}
	\newcommand{\examnum}{Exercício - Password} \newcommand{\examdate}{$\today$}
	\newcommand*{\oldneg}{\mathord{\nao}}
	\pagestyle{headandfoot}
	\runningheadrule
	\firstpageheader{}{}{}
	\runningheader{\class}{\examnum}{\examdate}
	\runningfootrule
	\firstpagefooter{}{}{}
	\runningfooter{}{Página \thepage\ de \numpages}{}
	%\pagestyle{foot}
	%\lfoot{}
	%\cfoot[]{Página \thepage\ de \numpages}
	%\rfoot{}
	%
	%\runningfootrule
	%\boxedpoints %\pointsinmargin
	%\printanswers
	%\SolutionEmphasis{\color{red}}
	%\definecolor{SolutionBoxColor}{gray}{0.8}
	%\renewcommand{\solutiontitle}{\noindent\textbf{Solução:}\enspace}
	%\renewcommand{\solutiontitle}{\noindent\textbf{Solution:}\par\noindent}

	\begin{document}

	\noindent
	\begin{tabular*}{\textwidth}{l @{\extracolsep{\fill}} r @{\extracolsep{6pt}} l}
	\textbf{Disciplina:} \class &  \multicolumn{2}{c}{\multirow{4}{*}{
	\includegraphics[scale=0.35]{logo.png}}}\\
	\textbf{Prof.:} ~\profname \vspace{0.2in} && \\
	\vspace{0.2in} \textbf{Nome:} \makebox[3in]{\hrulefill} \vspace{0.2in} &&\\
	%\textbf{Matrícula:} \makebox[1.5in]{\hrulefill} \textbf{Nota:} \makebox[1in]{\hrulefill} &&\\ 
	%\multicolumn{3}{c}{\textbf{{\examnum \hspace{3px} \examdate}}}\\
	\multicolumn{3}{c}{}\\
	\end{tabular*}\\
	\rule[2ex]{\textwidth}{1.6pt}
	\parskip=0.1in

	%Este exame contêm \numpages\ páginas (incluindo esta capa) e \numquestions\ questões.\\
	%Total de \numpoints~ pontos.

	%\begin{center}
	%Grade Table (for teacher use only)\\
	%\addpoints
	%\gradetable[v][questions]
	%\end{center}

	\noindent
\begin{questions}

% ===== FILOSOFIA E HISTÓRIA =====

\question Em 1970, no AI Lab do MIT, Richard Stallman recusava-se a usar senhas 
e definia a sua como ``Enter'' (senha vazia), para que qualquer pessoa pudesse 
acessar a sua conta. Ele justificava: ``não acredito que a segurança num 
computador seja realmente desejável''.

\begin{itemize}
\item Você concorda ou discorda da posição de Stallman? Justifique.
\item O contexto dele era um computador \textbf{compartilhado} numa comunidade 
de pesquisa. Isso muda a sua avaliação? Por quê?
\item Você acha que algo parecido seria possível hoje? Em que contexto?
\end{itemize}
\begin{solutionbox}{2in}
\end{solutionbox}

% ===== PANORAMA PASSWORDLESS =====

\question O que está substituindo as senhas? Cite e explique brevemente 
os quatro principais métodos: \textbf{Passkeys}, \textbf{Biometria}, 
\textbf{Chaves de segurança físicas} e \textbf{Autenticação comportamental}.
\begin{solutionbox}{3in}
\end{solutionbox}

\question Compare os quatro métodos em termos de \textbf{modelo de confiança} 
e \textbf{resistência a phishing}. Qual deles oferece a proteção mais forte 
contra phishing? Por quê?
\begin{solutionbox}{2.3in}
\end{solutionbox}

\question Explique o que é o padrão \textbf{FIDO2/WebAuthn}. Como ele se 
relaciona com Passkeys? Por que ele é considerado um avanço em relação às 
senhas tradicionais?
\begin{solutionbox}{2.5in}
\end{solutionbox}

% ===== CONCEITO-CHAVE =====

\question Explique o conceito de \textbf{origin binding} em Passkeys. 
Por que ele torna o Passkey resistente a phishing de forma \textbf{estrutural} 
(e não apenas ``mais difícil de enganar'')?
\begin{solutionbox}{2.5in}
\end{solutionbox}

% ===== RECUPERAÇÃO DE CONTA =====

\question Um sistema oferece Passkey como login principal, mas permite 
``recuperação de conta'' via Magic Link enviado por e-mail. Um atacante 
consegue comprometer a conta de e-mail do usuário.

\begin{itemize}
\item Qual é a estratégia de ataque mais direta que o atacante pode adotar?
\item Como esse cenário se relaciona com o lab de ``Password reset broken logic''?
\end{itemize}
\begin{solutionbox}{3.2in}
\end{solutionbox}

\question A \textbf{recuperação de conta} é um dos pontos mais críticos em 
segurança. Analise os cenários:

\begin{itemize}
\item Um usuário morre subitamente e a família precisa acessar as contas dele.
\item Um usuário usa apenas Passkey e perde o dispositivo.
\item Um usuário perde o acesso ao e-mail de recuperação.
\end{itemize}

Para cada cenário, discuta os riscos e o trade-off entre \textbf{segurança} 
e \textbf{recuperabilidade}.
\begin{solutionbox}{3.2in}
\end{solutionbox}

% ===== CRIPTOMOEDAS =====

\question Criptomoedas não usam ``senhas'' no sentido tradicional. Elas usam 
\textbf{chaves privadas} e \textbf{seed phrases}.

\begin{itemize}
\item Qual é a diferença fundamental entre uma senha de site e uma seed phrase?
\item Por que perder a seed phrase significa perder os fundos permanentemente?
\item Que estratégias você recomendaria para proteger uma seed phrase?
\item Quem é a ``fonte da verdade'' numa carteira de criptomoedas?
\end{itemize}
\begin{solutionbox}{5in}
\end{solutionbox}

\end{questions}

\end{document}
